\documentclass[10pt,a4paper]{amsart}
\usepackage[margin=19mm]{geometry}
\usepackage{amsmath,amssymb,mathtools}
\usepackage[scr=rsfs,cal=euler]{mathalfa}
\usepackage{xfrac}
\usepackage[unicode,hidelinks,bookmarks=false]{hyperref}
\newcommand{\ie}{i.e.\ }
\newcommand{\cf}{cf.\ }
\newcommand{\resp}{respectively\ }
\newcommand{\Subsection}{\subsection}
\providecommand{\nobreakdash}{\nobreak\hbox{-}}
\setlength{\emergencystretch}{2em}
\multlinegap0pt
\usepackage{amssymb}
\usepackage{xcolor}
\usepackage{algorithm}\usepackage{algpseudocode}
\usepackage{mathtools}
\usepackage{tikz}
\newtheorem{lemma}{Lemma}
\title{M-duality in fixed-point problems for general normed spaces}
\author{Jongmin Lee, Sangyeop Woo, Ernest K. Ryu}
\date{Source: arXiv:2609.08151v1 (CC BY 4.0)}
\begin{document}
\maketitle

\noindent\emph{Source and licence.} M-duality in fixed-point problems for general normed spaces. Version arXiv:2609.08151v1 (CC BY 4.0), distributed by arXiv under the Creative Commons Attribution 4.0 International licence, http://creativecommons.org/licenses/by/4.0/, as recorded in the arXiv licence metadata for this exact version and shown on the abstract page https://arxiv.org/abs/2609.08151v1. Full licence notice: \texttt{licenses/arxiv-2609.08151-CC-BY-4.0.txt}.\par\smallskip

\noindent\textit{Editorial scope.} Counted unit: the article's lemma lemma (source0.tex lines 1150--1160). The article's complete original proof of it is delivered with it (source0.tex lines 1161--1239, one \texttt{proof} environment of 79 lines). No mathematics is written, repaired or re-derived: the statement and the proof are the article's own lines, in the article's own notation, and no retained line is substituted or re-flowed.  No further source-local context is needed: every cross-reference of every retained line resolves inside the two delivered spans.  Nothing was omitted from any retained span, so the package carries no omission record.  No retained line contains a citation, so the wrapper carries no bibliography and no bibliography entry.  No retained line contains a figure environment, an image-inclusion command or drawing code, so no image is delivered and the package declares no asset dependency.  Editorial formatting only, in addition: an AMS \texttt{amsart} preamble that restates the article's own declarations and macros used by the retained lines, namely the environments \texttt{lemma} on separate counters, together with the standard packages those lines need.  The retained environments print in this document as Lemma 1 in order of appearance, whereas the article prints the same items with the numbering of its own full document.  The complete native source of the article as distributed in the arXiv e-print for this exact version is bundled unchanged as \texttt{sources/209699/source0.tex}.
\begin{lemma}
Let \(5/8<\ell\le 1\) and \(N\ge 2\), and assume that \(m^N<1\). Then
% \[
% g_N-\frac{1}{4(N+2)}
% >0
% \]
\[
g_N-\frac{1-m^N}{8}
\ge 0.
\]
\end{lemma}

\begin{proof}
Since $\frac{1+x_{n-1}^2}{2}$ with $x_0=0$ increases with $n$, by substituting $x_n:=\ell m^n < \ell $, we have
\[
x_n
=\frac{1+x_{n-1}^2}{2}.
\]
% In particular,
% \[
% x_1=\frac12,
% \qquad
% x_2=\frac58,
% \qquad
% 0<x_n<1.
% \]

Let
\[
X_n:=\prod_{r=1}^n x_r,
\qquad
P_n:=\frac{X_n}{1-x_n}.
\]
Since
\[
1-x_{n+1}
=
1-\frac{1+x_n^2}{2}
=
\frac{(1-x_n)(1+x_n)}{2},
\]
we have, for \(2\le n<N\),
\[
\begin{aligned}
\frac{P_{n+1}}{P_n}
=
\frac{x_{n+1}(1-x_n)}{1-x_{n+1}} =
\frac{1+x_n^2}{1+x_n}
<1,
\end{aligned}
\]
where the last inequality follows from \(0<x_n<\ell\).

Moreover, since
\[
P_2
=
\frac{x_1x_2}{1-x_2}
=
\frac{\frac12\cdot\frac58}{1-\frac58}
=
\frac56,
\]
we have $P_N \le \frac{5}{6}$. Also,
\[
g_N
=
1-x_N-X_N
=
(1-x_N)(1-P_N)
\]
 implies
\[
\begin{aligned}
g_N-\frac{1-m^N}{6}
&=
(1-x_N)(1-P_N)-\frac{1-m^N}{6} \\
&=
(1-x_N)\left(\frac56-P_N\right)
 +\frac{1-x_N-(1-m^N)}{6} \\
&=
(1-x_N)\left(\frac56-P_N\right)
 +\frac{(1-\ell)m^N}{6}.
\end{aligned}
\]
Both terms on the right-hand side are nonnegative. Hence
\[
g_N\ge \frac{1-m^N}{6}.
\]
This proves the claim.
\end{proof}

\end{document}
